% ---------------------------------------------------------------
% Chapter 7: Results
% ---------------------------------------------------------------
\chapter{Results} \label{cap:results}

\section{Tests}

Correctness is verified by three complementary test strategies: \textit{ISA compliance} tests, which
statically check that the compiler never emits a forbidden instruction; \textit{behavioral
self-tests}, which execute compiled programs on a reference simulator and check their results; and
the \textit{GCC torture suite}, which subjects the compiler to a large corpus of programs accumulated
by the GCC project itself. The synthesis-heavy targets are verified in depth: rvsc1 receives all
three layers, and rvsc0 receives the first two (the torture suite requires function calls, which
rvsc0 cannot compile). rvsc2 is covered by ISA compliance alone, and rvsc3 through rvsc7 require no
testing beyond what the upstream backend already provides.

\subsection{Test Strategies}

\subsubsection{ISA Compliance} \label{sec:strategy-isa}

The ISA compliance strategy answers a static question: does the compiler ever emit an instruction that
the modeled processor does not implement? Each target has a corpus of small C programs, each written
to exercise one operation category. The test script compiles every program in the corpus at each of
GCC's five optimization levels (\texttt{-O0}, \texttt{-O1}, \texttt{-O2}, \texttt{-O3},
\texttt{-Os}), assembles the output, and disassembles the resulting object with pseudo-instruction
expansion enabled. Every mnemonic in the disassembly is then checked against the target's allowlist.

\subsubsection{Behavioral Self-Tests} \label{sec:strategy-behav}

ISA compliance proves that the output is \textit{legal}, it says nothing about whether it is
\textit{correct}. The behavioral strategy closes this gap by executing the compiled programs on Spike,
the RISC-V reference simulator \cite{spike}, and checking their results. Each behavioral test is a
self-validating C program covering one operation category through a sequence of assertions in boundary
values, sign transitions, every byte lane of a word, and algebraic identities cross-checked against
independently computed results. The program returns 0 when every assertion holds, and a distinct
nonzero code identifying the first failing assertion otherwise.

\subsubsection{GCC Torture Suite} \label{sec:strategy-torture}

Hand-written tests only cover the code paths their author anticipated. The third strategy probes the
remaining ones with \texttt{gcc.c-torture/execute}, a suite of 1\,684 C programs accumulated by the
GCC project over three decades of compiler development, many distilled from real miscompilation bugs.
The programs are self-validating, with assertion calling \texttt{abort()} on
failure. The harness compiles each program at all five
optimization levels and runs it on Spike, giving $1684 \times 5 = 8420$ compiler/optimizer
combinations. A combination passes when the simulator exits with code 0.

\subsection{rvsc0}

\subsubsection{ISA Compliance} \label{sec:sc0-isa-tests}

The rvsc0 corpus contains 12 programs. Every mnemonic they produce is checked against the
eight-instruction rvsc0 allowlist:

\begin{lstlisting}
lw  sw  beq  add  addi  sub  and  or
\end{lstlisting}

\begin{table}[htb]
\centering
\caption{ISA compliance test programs for rvsc0, by operation category}
\label{tab:sc0-isa-files}
\begin{tabular}{ll}
\toprule
\textbf{Test program} & \textbf{Operations exercised} \\
\midrule
Arithmetic    & ADD, ADDI, SUB (regression: must remain native) \\
Branches      & BNE, BLT, BGE, BLTU, BGEU --- synthesized from \texttt{beq} via SLT chains \\
Signed byte load & LB --- synthesized via \texttt{lw}+shift+sign-extend \\
Halfword load & LH signed/unsigned --- synthesized via \texttt{lw}+shift+mask \\
Logic         & AND, OR (native); XOR via \texttt{(a|b)-(a\&b)}; ANDI/ORI via \texttt{li}+register-op \\
Large constants & LUI --- synthesized via constant pool: \texttt{lw rd, \%lo(pool)(x0)} \\
Bitwise NOT   & NOT --- synthesized as \texttt{sub x0, rs; addi rd, rd, -1} \\
Byte store    & SB --- synthesized via \texttt{lw}+clear+insert+\texttt{sw} \\
Shifts        & SLL, SRL, SRA with constant and variable shift counts \\
Comparisons   & SLT and SLTU --- synthesized via sub/xor/and/lshr chains \\
Variable arithmetic right shift & SRA with runtime shift count --- bit-extraction loop + sign-fill \\
Variable logical right shift & SRL with runtime shift count --- bit-extraction loop \\
\bottomrule
\end{tabular}
\end{table}

Each program is compiled at five optimization levels, giving $12 \times 5 = \mathbf{60}$ test cases
in total. All 60 pass: no forbidden mnemonic appears in any rvsc0 output at any optimization level.

\subsubsection{Behavioral Self-Tests} \label{sec:sc0-behav-tests}

Because rvsc0 has no \texttt{jalr} instruction, it cannot use the proxy-kernel runtime that rvsc1
uses. Instead, each test program is a single C function, linked against a small hand-written
bare-metal startup routine that sets up the stack, invokes the test function, converts its return
value to a host-interface exit token, and writes it to the simulator's designated exit address. Spike
runs the binary bare-metal at its default load address of \texttt{0x80000000} and exits with the
reported value. The test passes if Spike exits with code 0.

// TODO: fix it, I first assume the program always start on 0x but this is only true whn the host has virtual memory
A key constraint distinguishes rvsc0 behavioral tests from rvsc1: global variables and large integer
constants are forbidden. The rvsc0 constant pool is valid only when pool entries resolve to addresses
below 2048 (the 12-bit signed offset range of \texttt{x0}). At Spike's load address of
\texttt{0x80000000}, pool entries would be accessed via \texttt{lw rd, \%lo(pool)(x0)} with a
wrapped address, producing incorrect values. All test programs therefore use only stack-allocated
\texttt{volatile} locals and constants within the SMALL\_OPERAND range ($-2048$ to $2047$).

\begin{table}[htb]
\centering
\caption{Behavioral test programs for rvsc0, by operation category}
\label{tab:sc0-behav-files}
\begin{tabular}{ll}
\toprule
\textbf{Test program} & \textbf{Cases covered} \\
\midrule
Arithmetic  & ADD/SUB/ADDI on positive, negative, and zero operands; AND and OR identity and absorption \\
Branches    & Signed and unsigned comparisons: \texttt{<}, \texttt{>}, \texttt{<=}, \texttt{>=}, \texttt{!=}; all synthesized from \texttt{beq}+SLT chains \\
Logic       & XOR, ANDI, ORI on representative values; \texttt{(a|b)-(a\&b)} identity; complement-via-XOR \\
Loops       & Ascending for-loop (sum 1..10), countdown while-loop (doubling to 256), do-while (repeated addition), nested loops \\
Memory      & SB/LBU/LB on all four byte lanes; SH/LHU/LH on both halfword lanes; signed widening via volatile intermediary \\
Bitwise NOT & NOT on 0, $-1$, 1, $-128$, 127; combined \texttt{\~{}\&}, \texttt{\~{}|}; XOR cross-checked \\
Shifts      & SLL/SRL/SRA with constant counts (1, 3, 8); variable counts; sign-propagation (SRA) and zero-fill (SRL) \\
Comparisons & SLT and SLTU: signed ordering, unsigned ordering, equality; unsigned wrap-around larger than small positive \\
\bottomrule
\end{tabular}
\end{table}

All 8 behavioral tests pass at every optimization level (\texttt{-O0} through \texttt{-Os}), for
$8 \times 5 = 40$ cases.

\subsection{rvsc1}

\subsubsection{ISA Compliance} \label{sec:sc1-isa-tests}

The rvsc1 corpus contains 19 programs. Every mnemonic they produce is checked against the
ten-instruction rvsc1 allowlist:

\begin{lstlisting}
lw  sw  beq  add  addi  sub  and  or  lui  jalr
\end{lstlisting}

\begin{table}[htb]
\centering
\caption{ISA compliance test programs for rvsc1, by operation category}
\label{tab:sc1-isa-files}
\begin{tabular}{ll}
\toprule
\textbf{Test program} & \textbf{Operations exercised} \\
\midrule
Addition     & ADD and ADDI (regression: must remain native) \\
Immediate AND & ANDI --- synthesized as \texttt{li t, imm; and rd, rs, t} \\
Branches     & BNE, BLT, BGE, BLTU, BGEU --- all synthesized from \texttt{beq} \\
Function call & Call and return --- JAL synthesized as \texttt{lui+addi+jalr} \\
Signed byte load & LB --- synthesized via \texttt{lw}+shift+sign-extend \\
Unsigned byte load & LBU --- synthesized via \texttt{lw}+shift+mask \\
Signed halfword load & LH --- synthesized via \texttt{lw}+shift+sign-extend \\
Unsigned halfword load & LHU --- synthesized via \texttt{lw}+shift+mask \\
Loop         & Loop with synthesized branch and variable shift \\
Bitwise NOT  & NOT --- synthesized as \texttt{sub x0, rs; addi rd, rd, -1} \\
Immediate OR & ORI --- synthesized as \texttt{li t, imm; or rd, rs, t} \\
Byte store   & SB --- synthesized via \texttt{lw}+clear+insert+\texttt{sw} \\
Halfword store & SH --- synthesized via \texttt{lw}+clear+insert+\texttt{sw} \\
Constant-count shifts & SLL, SRL, SRA with compile-time shift counts \\
Variable left shift & SLL with runtime shift count --- count-down loop \\
Comparisons  & SLT and SLTU --- synthesized via sub/xor/and/lshr \\
Variable arithmetic right shift & SRA with runtime shift count --- bit-extraction loop + sign-fill \\
Variable logical right shift & SRL with runtime shift count --- bit-extraction loop \\
XOR          & XOR --- synthesized via \texttt{(a|b) - (a\&b)} \\
\bottomrule
\end{tabular}
\end{table}

Each program is compiled at five optimization levels, giving $19 \times 5 = \mathbf{95}$ test cases
in total. All 95 pass: no forbidden mnemonic appears in any rvsc1 output at any optimization level.

\subsubsection{Behavioral Self-Tests} \label{sec:sc1-behav-tests}

rvsc1 binaries run on top of the RISC-V proxy kernel, a minimal execution environment that loads the
program, provides a stack, and services the \texttt{exit} system call, allowing test programs to be
ordinary C programs with a \texttt{main} function.

\begin{table}[htb]
\centering
\caption{Behavioral test programs for rvsc1, by operation category}
\label{tab:sc1-behav-files}
\begin{tabular}{ll}
\toprule
\textbf{Test program} & \textbf{Cases covered} \\
\midrule
Branches    & Signed and unsigned comparisons: \texttt{<}, \texttt{>}, \texttt{<=}, \texttt{>=}, \texttt{!=}; boundary values including \texttt{INT\_MIN}, \texttt{INT\_MAX}, and \texttt{UINT\_MAX} \\
Function calls & Recursive Fibonacci (\texttt{fib(10)=55}), multi-argument calls, call through a function pointer \\
Logic       & XOR, ORI, ANDI, NOT on representative bit patterns; \texttt{(a|b)-(a\&b)} identity \\
Memory      & SB/LBU/LB and SH/LHU/LH on aligned addresses; signed/unsigned byte and halfword widening \\
Shifts      & SLL/SRL/SRA with constant counts (3, 15, 16, 4); variable counts; sign-propagation and zero-fill \\
Comparisons & SLT and SLTU with \texttt{a=-1}, \texttt{b=1}; \texttt{INT\_MIN < INT\_MAX}; \texttt{UINT\_MAX > 0} \\
\bottomrule
\end{tabular}
\end{table}

All 6 behavioral tests pass at every optimization level (\texttt{-O0} through \texttt{-Os}), for
$6 \times 5 = 30$ cases.

\subsubsection{GCC Torture Suite} \label{sec:sc1-torture-tests}

The full suite covers $1684 \times 5 = 8420$ compiler+optimizer combinations.
Table~\ref{tab:torture-results} summarizes the outcome.

\begin{table}[htb]
\centering
\caption{gcc.c-torture/execute results for rvsc1 (8\,420 combinations)}
\label{tab:torture-results}
\begin{tabular}{lrrr}
\toprule
\textbf{Optimization} & \textbf{Passed} & \textbf{Skipped} & \textbf{Timed out} \\
\midrule
\texttt{-O0} & $\sim$1\,501 & $\sim$183 & 0 \\
\texttt{-O1} & $\sim$1\,531 & $\sim$153 & 0 \\
\texttt{-O2} & 1\,505  & 164  & 15 \\
\texttt{-O3} & $\sim$1\,503 & $\sim$161 & $\sim$14 \\
\texttt{-Os} & $\sim$1\,525 & $\sim$157 & $\sim$2 \\
\textbf{Total} & \textbf{7\,546} & \textbf{843} & \textbf{31} \\
\bottomrule
\end{tabular}
\end{table}

``Skipped'' denotes programs that fail to compile or link. The two dominant causes are unrelated to
sc1's restricted instruction set: roughly 100 tests use old K\&R-style implicit \texttt{int}
declarations that GCC 17 (defaulting to C23) rejects as hard errors, and roughly 35 tests use
\texttt{printf}/\texttt{sprintf}/\texttt{fprintf}, whose newlib implementation references an internal
reentrant symbol absent from the linked sysroot. A further three use \texttt{\_\_int128}, which
32-bit RISC-V does not support regardless of instruction set.

The 31 timed-out cases are programs that compile and execute correctly but generate so many
synthesized instructions that Spike exceeds the 300-second simulation budget.
Table~\ref{tab:torture-timeouts} lists the programs that consistently time out at \texttt{-O2} or
\texttt{-O3}.

\begin{table}[htb]
\centering
\caption{gcc.c-torture programs that consistently time out ($>$300\,s on Spike) due to synthesis overhead}
\label{tab:torture-timeouts}
\begin{tabular}{ll}
\toprule
\textbf{Test program} & \textbf{Optimization levels} \\
\midrule
\texttt{920302-1.c}   & \texttt{-O2}, \texttt{-O3} \\
\texttt{pr50865.c}    & \texttt{-O2}, \texttt{-O3} \\
\texttt{pr53645-2.c}  & \texttt{-O2}, \texttt{-O3} \\
\texttt{pr57876.c}    & \texttt{-O3} \\
\texttt{pr69447.c}    & \texttt{-O2}, \texttt{-O3} \\
\texttt{pr82524.c}    & \texttt{-O2}, \texttt{-O3} \\
\texttt{pr91450-1.c}  & \texttt{-O2}, \texttt{-O3} \\
\texttt{pr91450-2.c}  & \texttt{-O2}, \texttt{-O3} \\
\texttt{pr93249.c}    & \texttt{-O3} \\
\texttt{pr93908.c}    & \texttt{-O2}, \texttt{-O3} \\
\texttt{string-opt-5.c} & \texttt{-O2}, \texttt{-O3} \\
\bottomrule
\end{tabular}
\end{table}

\subsection{rvsc2} \label{sec:sc2-isa-tests}

rvsc2 removes only \texttt{fence} (and \texttt{fence.i}) from RV32I, and nothing is synthesized:
every instruction the compiler may emit is native. The only property left to verify is ISA
compliance. The rvsc1 test corpus is reused for this check, but against a much wider allowlist: the
full RV32I base set minus the fence, CSR, and system instruction groups. The 19 programs at five
optimization levels give 95 test cases; all pass.

\subsection{rvsc3 and Above}

No testing is performed for rvsc3 through rvsc7. These targets contain no synthesis code: they only
enable flags that correspond to instructions already present in the upstream RISC-V backend. Their
correctness follows from the correctness of the upstream backend and the GCC test suite.

\section{Program Size}

Each synthesized instruction expands into a sequence of native instructions, increasing the static
size of the compiled binary. Synthesis sequences fall into two categories. \textit{Constant-length}
expansions always emit the same number of instructions regardless of operand values: NOT expands to 2
instructions, XOR to 3 (register operands; 4 with an immediate operand), and each immediate variant
(ANDI, ORI) adds 1 instruction. \textit{Variable-length} expansions depend on runtime values: SLL,
SRL, and SRA use count-down loops whose length is proportional to the shift amount, with worst-case
counts of 187, $\sim$170, and $\sim$200 instructions respectively for a shift of 31.

To quantify the size penalty on realistic code, the Embench-IoT benchmark suite \cite{embench} was
compiled for three toolchains:

\begin{itemize}
  \item \textbf{gcc17} --- the same GCC 17 fork used by this work, built for the stock
        \texttt{riscv32-unknown-elf} triple with no rvsc flags (\texttt{-march=rv32i -mabi=ilp32}),
        serving as the upstream reference;
  \item \textbf{rvsc2} --- the native RV32I target of this work, in which nothing is synthesized;
        and
  \item \textbf{rvsc1} --- the synthesized target.
\end{itemize}

Code size is measured as the \texttt{.text} section of each benchmark's own object files. All
benchmarks were compiled at \texttt{-O2}.

\begin{table}[htb]
\centering
\caption{Embench-IoT static code size (\texttt{.text} bytes, \texttt{-O2}) and rvsc1 expansion ratios}
\label{tab:embench-size}
\begin{tabular}{lrrr}
\toprule
\textbf{Benchmark} & \textbf{gcc17} & \textbf{rvsc1} & \textbf{rvsc1 / gcc17} \\
\midrule
\texttt{aha-mont64}     &  3\,548 &  17\,940 &  5.06$\times$ \\
\texttt{crc32}          &    380  &     908  &  2.39$\times$ \\
\texttt{depthconv}      &    608  &   3\,532 &  5.81$\times$ \\
\texttt{edn}            &  3\,244 &  22\,832 &  7.04$\times$ \\
\texttt{huffbench}      &  2\,140 &   8\,688 &  4.06$\times$ \\
\texttt{matmult-int}    &    936  &   1\,476 &  1.58$\times$ \\
\texttt{md5sum}         &  1\,016 &   2\,408 &  2.37$\times$ \\
\texttt{nettle-aes}     &  4\,444 &  46\,872 & 10.55$\times$ \\
\texttt{nettle-sha256}  &  6\,976 &  44\,588 &  6.39$\times$ \\
\texttt{nsichneu}       & 19\,668 &  65\,364 &  3.32$\times$ \\
\texttt{picojpeg}       & 15\,360 & 149\,496 &  9.73$\times$ \\
\texttt{qrduino}        & 12\,852 & 126\,048 &  9.81$\times$ \\
\texttt{sglib-combined} & 10\,824 &  58\,256 &  5.38$\times$ \\
\texttt{slre}           &  4\,256 &  27\,720 &  6.51$\times$ \\
\texttt{statemate}      &  6\,484 & 106\,244 & 16.39$\times$ \\
\texttt{tarfind}        &    528  &   2\,620 &  4.96$\times$ \\
\texttt{ud}             &  1\,436 &   2\,020 &  1.41$\times$ \\
\texttt{wikisort}       &  7\,744 &  20\,032 &  2.59$\times$ \\
\texttt{xgboost}        &    624  &   5\,620 &  9.01$\times$ \\
\bottomrule
\end{tabular}
\end{table}

Across the suite, synthesis inflates code size by a mean of $4.94\times$ relative to the
\texttt{gcc17} baseline. The custom \texttt{rvsc2} target produces code identical to \texttt{gcc17}
for all 19 benchmarks (geomean ratio = 1.000).

\section{Program Performance}

Since the target processor is single-cycle, every instruction retires in exactly one clock cycle
(ignoring memory latency). Dynamic instruction counts are obtained by executing each compiled binary
on Spike and totalling the retired instructions.

\subsection{Whole-Program Instruction Counts} \label{sec:embench-perf}

The dynamic cost of a whole program depends on how often each synthesized instruction executes 
at run time, simple instructions such not or xori has a small absolute value, other instructions
has a loop and the cost change depends on params. This brenchmark tries to measure a real program impact  weighted by loop trip counts rather than by static frequency. 

Five benchmarks exceeded a 300-second Spike budget under rvsc1 and are reported as TIMEOUT.

\begin{table}[htb]
\centering
\caption{Embench-IoT dynamic retired-instruction counts on Spike and rvsc1 run-time overhead.
         Geomean is over the thirteen benchmarks that completed on both toolchains.}
\label{tab:embench-perf}
\begin{tabular}{lrrr}
\toprule
\textbf{Benchmark} & \textbf{rvsc2 (native)} & \textbf{rvsc1 (synth)} & \textbf{Overhead} \\
\midrule
\texttt{aha-mont64}     & 12\,997\,428 & \textit{timeout}    & --- \\
\texttt{crc32}          &  5\,991\,188 & 130\,556\,929  & 21.8$\times$ \\
\texttt{depthconv}      & 54\,552\,784 & 1\,306\,486\,627 & 23.9$\times$ \\
\texttt{edn}            & 68\,624\,580 & 512\,612\,605  &  7.5$\times$ \\
\texttt{huffbench}      &  2\,399\,142 & \textit{timeout}    & --- \\
\texttt{matmult-int}    & 24\,445\,873 &  25\,543\,407  &  1.0$\times$ \\
\texttt{md5sum}         &  2\,917\,425 &  38\,694\,462  & 13.3$\times$ \\
\texttt{nettle-aes}     &  4\,674\,471 & 171\,446\,978  & 36.7$\times$ \\
\texttt{nettle-sha256}  &  4\,917\,841 & \textit{timeout}    & --- \\
\texttt{nsichneu}       &  2\,514\,597 &  14\,689\,559  &  5.8$\times$ \\
\texttt{picojpeg}       &  3\,602\,866 & \textit{timeout}    & --- \\
\texttt{qrduino}        &  5\,140\,049 & \textit{timeout}    & --- \\
\texttt{sglib-combined} &  3\,005\,822 &  68\,953\,363  & 22.9$\times$ \\
\texttt{slre}           &  2\,983\,415 &  88\,220\,929  & 29.6$\times$ \\
\texttt{statemate}      &  2\,065\,897 & 150\,934\,719  & 73.1$\times$ \\
\texttt{tarfind}        &  5\,287\,906 &  37\,627\,640  &  7.1$\times$ \\
\texttt{ud}             &  6\,622\,675 &   8\,378\,022  &  1.3$\times$ \\
\texttt{wikisort}       & \textit{linkfail} & 17\,710\,935 & --- \\
\texttt{xgboost}        &  3\,804\,604 & 433\,520\,452  & 113.9$\times$ \\
\bottomrule
\end{tabular}
\end{table}

Across the thirteen benchmarks that completed on both toolchains, synthesis inflates the dynamic
instruction count by a mean of $13.7\times$, roughly triple the $4.94\times$ static-size
penalty, because the most expensive syntheses sit inside the hottest loops.

\section{Discussion}

The results establish correctness first and cost second. On the correctness axis, the ISA compliance
tests confirm that the compiler never emits a forbidden mnemonic. Behavioral equivalence is established independently by differential execution on Spike: every
rvsc1 program produces the same exit code as the reference RV32I binary compiled from the same
source, and every rvsc0 single-function program writes the same \texttt{tohost} value as its
reference.

The cost of that re-expression is quantified along two dimensions. Statically, synthesis inflates
code size by a mean of $4.94\times$ over the Embench suite
(Table~\ref{tab:embench-size}), dynamically, it inflates the retired-instruction count by a
mean of $13.7\times$ over the benchmarks that complete (Table~\ref{tab:embench-perf}). The dynamic
penalty is the larger of the two because the costliest syntheses of the shift loops, each of which
re-materializes its own back-edge every iteration (Section~\ref{sec:sc1-sll}), tend to sit inside
the hottest loops, so their cost is multiplied by trip count rather than merely by static occurrence.

For the pedagogical setting these targets are built for, this variation is the point rather than a
limitation. The programs students write in an introductory single-cycle course, small loops, modest
shift amounts, few byte-granular memory accesses, fall at the inexpensive end of both distributions.
At the same time, the wide spread makes the cost of each ISA restriction concrete and measurable: a
student can compile the same source for rvsc1 and rvsc3, compare the \texttt{-S} output, and see
exactly how many native instructions a single missing \texttt{sll} or \texttt{sb} expands into.


% ---------------------------------------------------------------
% Chapter 8: Conclusion
% ---------------------------------------------------------------
\chapter{Conclusion} \label{cap:conclusion}

This work developed eight GCC compiler targets for the simplified RISC-V processors described in the
Hennessy--Patterson textbook, addressing the practical barrier that prevented students of PCS3225 at
USP from compiling and running C programs on the processors they build. The four specific objectives
stated in Chapter~\ref{cap:intro} were met: a minimal target (\texttt{rvsc0}) for the
eight-instruction Chapter 4.4 processor, a target (\texttt{rvsc1}) matching the course homework
extension with full C calling convention support, and six progressive targets (\texttt{rvsc2} through
\texttt{rvsc7}) covering the full RV32I-to-RV64IMAFD progression. Every instruction not natively
supported by a given target is synthesized from the primitives that target does provide, transparently
to the programmer.

\section{Contributions}

The primary contribution is the set of synthesis techniques embedded in the GCC machine description.
For the two most restricted targets, rvsc0 and rvsc1, eighteen distinct operations require synthesis,
ranging from one-instruction replacements (NOT, immediate variants) to variable-length loops (SLL,
SRL, SRA), multi-instruction identities (XOR, SLT, SLTU), read-modify-write sequences (LB, LBU, LH,
LHU, SB, SH), and call-site code generation (JAL, JMP). The rvsc0 target additionally requires
constant pool materialization for LUI, since 32-bit constants cannot otherwise be constructed from
the eight available instructions. Each synthesis was derived algebraically and embedded as a
\texttt{define\_expand} in \texttt{riscv.md}, so GCC selects and schedules the sequence as part of
normal compilation with no programmer intervention.

A secondary contribution is the validation methodology. Two independent test layers were developed
and applied: an ISA compliance suite that disassembles every generated object with
\texttt{objdump -M no-aliases} and verifies that no forbidden mnemonic appears, and a behavioral
equivalence suite that executes rvsc1 and rvsc0 binaries on Spike and compares their outputs against
a reference RV32I build. Together these layers confirm that synthesis is both correct by construction
(no forbidden instruction is ever emitted) and correct by execution (the computed results are
indistinguishable from those of a full-ISA compiler).

\section{Results Summary}

Correctness was established for all synthesis cases in rvsc0 and rvsc1. All ISA compliance tests
pass, and the differential behavioral tests confirm semantic equivalence across all tested programs.
The rvsc2 target compiles the full Embench-IoT suite with code identical to the upstream GCC 17
baseline (geomean ratio = 1.000), confirming that the custom target configuration introduces no
overhead relative to a stock build.

The cost of synthesis was quantified on two axes. Statically, the synthesized rvsc1 target inflates
code size by a mean of $4.94\times$ over the nineteen Embench benchmarks, with a range
from $1.41\times$ (\texttt{ud}, few shifts) to $16.39\times$ (\texttt{statemate},
comparison-heavy control flow). Dynamically, it inflates retired instruction counts by a
mean of $13.7\times$ over the thirteen benchmarks that completed within the time budget, with five
benchmarks timing out entirely.

\section{Limitations}

Synthesis does not apply to targets rvsc2 through rvsc7 --- these targets expose the full upstream
RISC-V backend and require no new synthesis logic. Their correctness depends entirely on the upstream
GCC test suite.

The rvsc0 target, unlike rvsc1, cannot execute programs that call and return from functions, because
\texttt{jalr} is absent. The constant pool mechanism enables 32-bit constant loading and
unconditional jumps, but programs must be written as non-returning single functions. This restriction
matches the processor it targets, but it means the behavioral test harness that uses HTIF (which
requires \texttt{jalr} for the call to \texttt{main}) cannot be used for rvsc0; rvsc0 programs write
their result directly to \texttt{tohost} via \texttt{sw}.

The synthesized shift loops are functionally correct but dynamically expensive to the point of
impracticality for programs that shift inside hot loops. The synthesis is inherently sequential: the
GCC machine description expands a shift at compile time into a counted loop, which at run time
executes one iteration per bit position. There is no partial-hardware or table-driven alternative
within the instruction set these targets support.

\section{Future Work}

The most direct extension would be performance optimizations within the existing synthesis.
Constant-amount shifts are currently synthesized as loops, even though the shift amount is known at
compile time and could instead be unrolled statically, reducing the per-shift cost from $6b + 1$ to
$2b$ instructions without any new hardware. This would require a separate \texttt{define\_expand}
branch for \texttt{CONST\_INT\_P} shift counts.

On the hardware side, each synthesis in this work corresponds exactly to the cost of the missing
instruction in hardware. A student who extends the Chapter 4.4 processor with, for example, a barrel
shifter could recompile with \texttt{-mshift} enabled and the compiler would switch to native
\texttt{sll}/\texttt{srl}/\texttt{sra} automatically, making the hardware improvement immediately
observable in both binary size and execution time. Building this feedback loop into a course lab
is a natural next step.

Finally, the target configuration and synthesis infrastructure developed here could be packaged as a
course resource for PCS3225, including pre-built toolchain binaries, startup files, linker scripts,
and a simple Makefile or justfile that lets students go from a C source file to a Spike execution
with a single command.
